\documentclass{llncs}
\usepackage[utf8]{inputenc}
\usepackage[spanish]{babel}

\begin{document}

\title{Machete de Metodologías: Agile, Scrum y Lean Six Sigma}
\author{Tu Nombre}
\institute{Tu Institución}

\maketitle

\begin{abstract}
Este documento sirve como referencia rápida sobre los conceptos fundamentales de la metodología Agile, el marco de trabajo Scrum, y los niveles de certificación White Belt y Green Belt (Lean Six Sigma).
\end{abstract}

\section{Metodología Ágil (Agile)}
Agile es una filosofía o conjunto de principios para la gestión de proyectos y desarrollo de software.
\begin{itemize}
    \item \textbf{Enfoque:} Iterativo e incremental. Entregas tempranas y continuas.
    \item \textbf{Valores principales:} Prioriza a los individuos y sus interacciones sobre los procesos; el producto funcionando sobre la documentación extensiva; la colaboración con el cliente sobre la negociación contractual; y la adaptación al cambio sobre el seguimiento estricto de un plan.
\end{itemize}

\section{Scrum}
Scrum es el marco de trabajo (framework) Ágil más utilizado para gestionar el desarrollo de productos complejos.
\begin{itemize}
    \item \textbf{Roles:} 
    \begin{itemize}
        \item \emph{Product Owner:} Representa al negocio, prioriza las necesidades y gestiona el Product Backlog.
        \item \emph{Scrum Master:} Actúa como facilitador (servant-leader), elimina impedimentos y asegura que el equipo entienda y aplique Scrum.
        \item \emph{Developers:} El equipo autoorganizado que ejecuta el trabajo técnico.
    \end{itemize}
    \item \textbf{Artefactos:}
    \begin{itemize}
        \item \emph{Product Backlog:} Lista priorizada de todos los requerimientos del producto.
        \item \emph{Sprint Backlog:} Tareas y objetivos seleccionados para el ciclo de trabajo actual.
        \item \emph{Incremento:} El producto funcional y terminado al final de cada ciclo.
    \end{itemize}
    \item \textbf{Eventos:} Sprint (ciclo de trabajo de 1 a 4 semanas), Sprint Planning, Daily Scrum (reunión diaria de 15 min), Sprint Review y Sprint Retrospective.
\end{itemize}

\section{Lean Six Sigma: White y Green Belt}
Aunque a menudo se integran con entornos Ágiles, los niveles "Belt" provienen de Lean Six Sigma, una metodología enfocada en la reducción de variabilidad, eliminación de desperdicios y mejora continua de procesos.
\begin{itemize}
    \item \textbf{White Belt (Cinturón Blanco):} Es el nivel introductorio. Los White Belts comprenden los conceptos básicos, el vocabulario de la metodología y la estructura fundamental de resolución de problemas (ciclo DMAIC: Definir, Medir, Analizar, Mejorar, Controlar). Participan en proyectos de mejora aportando su experiencia diaria, pero no los lideran.
    \item \textbf{Green Belt (Cinturón Verde):} Es un nivel operativo intermedio. Los Green Belts lideran proyectos de mejora de procesos dentro de su área específica de trabajo. Dedican una parte de su jornada laboral a estas tareas, aplicando herramientas estadísticas para analizar datos y resolver problemas crónicos, generalmente bajo la mentoría de un Black Belt.
\end{itemize}

\end{document}
